% Source: https://artofproblemsolving.com/community/c3h1699530_introduction_to_geometry_exercise_711_part_c_assumption
% Stable identity: aops-topic-1699530-q01
Let $r$ be the reflection in line $m$. Since $m$ is the perpendicular bisector of $\overline{AB}$, $r$ swaps $A$ and $B$; likewise it swaps $C$ and $D$.

(a) \textbf{Not necessarily.} Any two segments that are both bisected perpendicularly by $m$ work, whatever their lengths; for example $AB=4.4$ and $CD=2.4$ in the figure.

(b) \textbf{True.} Both $\overline{AB}$ and $\overline{CD}$ are perpendicular to $m$, and two lines perpendicular to the same line are parallel. (If the two segments happen to lie on the same line perpendicular to $m$, they share the midpoint where that line meets $m$, and they are collinear rather than parallel in the strict sense.)

(c) \textbf{True, in every case.} The reflection $r$ maps $A\mapsto B$ and $C\mapsto D$, so it maps segment $AC$ onto segment $BD$. Reflections preserve distances, so $AC=BD$.

\begin{asy}
usepackage("amsmath");
size(9cm);
defaultpen(fontsize(10pt));
real off=7;
draw((0,-0.6)--(0,3.4),linewidth(0.8));label("$m$",(0,3.4),N);
pair A=(-2.2,2.6),B=(2.2,2.6),C=(-1.2,0.6),D=(1.2,0.6);
draw(A--B,linewidth(0.9));draw(C--D,linewidth(0.9));
draw(A--C,dashed+blue);draw(B--D,dashed+blue);
dot(A);dot(B);dot(C);dot(D);
label("$A$",A,W);label("$B$",B,E);label("$C$",C,W);label("$D$",D,E);
label("$A$, $C$ on the same side",(0,-1.3));
draw((off,-0.6)--(off,3.4),linewidth(0.8));label("$m$",(off,3.4),N);
pair A2=(off-2.2,2.6),B2=(off+2.2,2.6),C2=(off+1.2,0.6),D2=(off-1.2,0.6);
draw(A2--B2,linewidth(0.9));draw(C2--D2,linewidth(0.9));
draw(A2--C2,dashed+blue);draw(B2--D2,dashed+blue);
dot(A2);dot(B2);dot(C2);dot(D2);
label("$A$",A2,W);label("$B$",B2,E);label("$C$",C2,E);label("$D$",D2,W);
label("$A$, $C$ on opposite sides",(off,-1.3));
label("\mbox{}",(min(currentpicture,user=true)+max(currentpicture,user=true))/2,black); // keeps the text after the figure black when LaTeX typesets the labels inline
\end{asy}

It does \emph{not} matter on which sides of $m$ the points $A$ and $C$ lie: the argument never uses it. The figure shows both cases; in the first $ACDB$ is an isosceles trapezoid, in the second the segments $AC$ and $BD$ cross on $m$. So no extra assumption is needed. In summary: $\boxed{\text{(a) not necessarily; (b) and (c) always}}$.
